\section{Discussion}
\label{sec:discussion}

We studied neural posterior estimation as an amortized approach to Bayesian
dynamic borrowing: pretrain once on a broad simulation distribution, then obtain
posteriors for new current--external pairs in a single forward pass. The
simulation study isolates covariate shift, outcome drift, and their combination;
ADNI illustrates that cohort mismatch can pull inference away from the
concurrent target even when discrimination remains strong. The empirical scope
is a scalar current-study target---an additive study-effect parameter in the
Gaussian simulations and a cohort-level risk parameter in ADNI. Covariate
effects, treatment--covariate interactions, and other causal functionals require
additional evaluation; Appendix~\ref{app:covariate-effect} provides only a
limited illustrative check for one target-aware covariate-effect analysis. This
framing keeps the real-data example aligned with the estimand actually analyzed
rather than treating ADNI as a randomized treatment-effect study.

Posterior quality depends on the pretraining simulator, so mismatch patterns
outside its support cannot be handled reliably. The current hand-crafted
summaries expose auditable signals of covariate shift, outcome-mechanism drift,
overlap, and joint mismatch, but may miss richer patient-level structure.
Runnable code is included in the supplementary material; ADNI raw data require
separate data-use access. Future work will study richer paired-data
representations, such as patient-level encoders for irregular longitudinal
records, posterior-spread calibration in severe mismatch regimes, multi-source
borrowing, and broader evaluation across estimands and prospective
data-borrowing settings.
